% !TeX root = ../../Book5.tex
% 纲要标题：### 19.6 与组合表示论的联系
% 纲要位置：工作记录/计划纲要.md，第 4349 行。
% * A 型与 Schur 函数
% * 对称幂、外幂及一般 Schur 函子
% * 第7章与最高权理论的互相解释
% * 综合本章已证公式核对 \(\mathfrak{sl}_2\)、\(\mathfrak{sl}_3\) 与 A 型 Schur 函数；第20章 BGG 方法保留为另一结构解释，本章不依赖这一方法
% ---
% # 第四编　高级表示论与代数对称

\section{与组合表示论的联系}
\label{b5:sec:1906}

第七章已经从对称群与张量幂构造 Schur 函子；本章则从 Cartan 子代数和根构造最高权表示。到 \(A\) 型，两种构造给出同一批表示。本节在一般 Weyl 公式已经证明之后建立这个识别，用来解释分拆和最高权之间的关系。

\subsection{A 型与 Schur 函数}
\label{b5:subsec:1906:01}

设 \(\mathfrak g=\mathfrak{sl}_r\)，取自然模 \(V=\mathbb C^r\)。
在对角 Cartan 子代数上，记 \(\varepsilon_i\) 为第 \(i\) 个对角坐标，并有
\(\sum_i\varepsilon_i=0\)。正根为 \(\varepsilon_i-\varepsilon_j\)、\(i<j\)，Weyl 群为 \(S_r\)，通过置换坐标作用。

\begin{proposition}\label{b5:prop:type-a-schur-character}
设 \(r\ge2\) 为整数，\(\mathfrak g=\mathfrak{sl}_r(\mathbb C)\)，取对角 Cartan 子代数及上三角正根，相应基本权记为 \(\omega_1,\ldots,\omega_{r-1}\)。设 \(\nu=(\nu_1\ge\cdots\ge\nu_r\ge0)\) 为分拆，令
\[
\lambda_\nu=\sum_{i=1}^{r-1}(\nu_i-\nu_{i+1})\omega_i.
\]
记 \(L(\lambda_\nu)\) 为相应有限维简单最高权模，\(s_\nu\) 为 Schur 多项式。用形式变量 \(x_1,\ldots,x_r\) 记录对角权，并规定 \(x_1\cdots x_r=1\)，则
\[
\operatorname{ch}L(\lambda_\nu)
=\frac{\det(x_i^{\nu_j+r-j})_{i,j=1}^r}
{\det(x_i^{r-j})_{i,j=1}^r}
=s_\nu(x_1,\ldots,x_r).
\]
\end{proposition}
\begin{proof}
指数列 \(\nu\) 在 \(\mathfrak{sl}_r\) 权格中的简单余根取值为
\(\nu_i-\nu_{i+1}\)，所以代表 \(\lambda_\nu\)。
列 \((r-1,r-2,\ldots,0)\) 的相邻差都是一，故代表 \(\rho\)；
其与通常和为零的 \(\rho\) 坐标只差共同常数，在 \(\sum_i\varepsilon_i=0\) 后为同一权。
Weyl 群交错求和正是变量置换的行列式，所以 Weyl 特征标公式给出前一个商。
后一个等号就是\cref{b5:def:schur-function}的交错式定义。
\end{proof}

每个支配整权 \(\sum_i a_i\omega_i\) 都可由唯一一个满足 \(\nu_r=0\) 的分拆表示：
\[
\nu_i=a_i+a_{i+1}+\cdots+a_{r-1},\qquad \nu_r=0.
\]
若允许末行非零，则对所有行同时加一个非负整数不会改变 \(\lambda_\nu\)。这正对应一般线性群表示中的行列式扭曲，在 \(\mathfrak{sl}_r\) 上则变成平凡扭曲。

\subsection{对称幂、外幂与 Schur 函子}
\label{b5:subsec:1906:02}

\begin{theorem}\label{b5:thm:schur-highest-identification}
设 \(r\ge2\) 为整数，\(V=\mathbb C^r\) 为 \(\mathfrak{sl}_r(\mathbb C)\) 的自然模，取对角 Cartan 子代数及上三角正根，相应基本权记为 \(\omega_1,\ldots,\omega_{r-1}\)。设 \(\nu=(\nu_1\ge\cdots\ge\nu_r\ge0)\) 为至多 \(r\) 行的分拆，不足 \(r\) 行时补零，令 \(\lambda_\nu=\sum_{i=1}^{r-1}(\nu_i-\nu_{i+1})\omega_i\)。Schur 函子值
\(\mathbb S_\nu(V)\) 按张量作用成为复 \(\mathfrak{sl}_r\) 模时，与相应有限维简单最高权模 \(L(\lambda_\nu)\) 满足
\[
\mathbb S_\nu(V)\cong L(\lambda_\nu).
\]
特别地，对整数 \(m\ge0\) 及 \(1\le j<r\)，有
\[
\operatorname{Sym}^mV\cong L(m\omega_1),\qquad
\bigwedge^jV\cong L(\omega_j)\quad(1\le j<r),\qquad
\bigwedge^rV\cong L(0).
\]
\end{theorem}
\begin{proof}
在 \(V^{\otimes n}\) 上，\(x\in\mathfrak{sl}_r\) 按
\[
x(v_1\otimes\cdots\otimes v_n)
=\sum_{a=1}^n v_1\otimes\cdots\otimes xv_a\otimes\cdots\otimes v_n
\]
作用；它与因子置换交换，所以在
\(\mathbb S_\nu(V)=\operatorname{Hom}_{S_n}(S^\nu,V^{\otimes n})\)
上也有自然作用。对角 Cartan 权由各张量指标的出现次数记录；
第七章的\cref{b5:prop:schur-functor-character}因此给出其形式特征标
\(s_\nu(x)\)，再加上 \(x_1\cdots x_r=1\) 的关系。

由\cref{b5:prop:type-a-schur-character}，这个特征标等于
\(\operatorname{ch}L(\lambda_\nu)\)。两个模有限维，故 Weyl 完全可约性和
\cref{b5:prop:formal-character-operations}说明它们同构。
单行与单列 Schur 函子分别为对称幂与外幂，代入相邻行长差得到后面的具体识别。
\end{proof}

这里的证明使用的是两套已经独立完成的理论。第七章的 Schur--Weyl 对偶不依赖本章；本章的一般最高权分类与 Weyl 特征标公式也没有借用第七章的 \(A\) 型结果。

\subsection{对称函数与最高权理论}
\label{b5:subsec:1906:03}

对 \(\mathfrak{sl}_r\) 的两个 Schur 模，先在对称函数环中求乘积
\[
s_\lambda s_\mu=\sum_\nu c_{\lambda\mu}^{\nu}s_\nu.
\]
删去超过 \(r\) 行的分拆，再按相邻行长差将每个剩余 \(\nu\) 换成最高权，
就得到对应李代数模的张量分解。若不同分拆给出同一最高权，须合并重数；
在同一总次数中，它们若只差共同整数平移，实际上平移量只能为零。
对称幂与外幂的乘法则分别由已经证明的横条、竖条 Pieri 规则直接计算。

\begin{example}\label{b5:exm:sl3-adjoint-square}
对 \(r=3\)，分拆 \((2,1,0)\) 对应 \(\omega_1+\omega_2\)，即八维伴随模。
由\eqref{b5:eq:lr-square-21}，保留至多三行的项并换成最高权，得到
\[
\begin{aligned}
L(\omega_1+\omega_2)^{\otimes2}
\cong{}&
L(2\omega_1+2\omega_2)\oplus L(3\omega_1)\oplus L(3\omega_2)\\
&\oplus L(\omega_1+\omega_2)^{\oplus2}\oplus L(0).
\end{aligned}
\]
具体对应的分拆依次为 \((4,2,0)\)、\((4,1,1)\)、\((3,3,0)\)、
\((3,2,1)\)、\((2,2,2)\)。维数公式给出
\(27+10+10+2\cdot8+1=64\)，与 \(8^2\) 一致。这里重数二来自
Littlewood--Richardson 系数；零权重数则属于每个简单模内部，是另一项数据。
\end{example}

\subsection{低秩计算与 BGG 方法简介}
\label{b5:subsec:1906:04}

对 \(\mathfrak{sl}_2\)，分拆 \((m,0)\) 的 Schur 多项式在
\(x_1=z,x_2=z^{-1}\) 时为
\(z^m+z^{m-2}+\cdots+z^{-m}\)，与本章秩一情形的 Weyl 计算一致。
对 \(\mathfrak{sl}_3\)，分拆 \((a+b,b,0)\) 正给出
\eqref{b5:eq:sl3-weyl-determinant}；三个变量的 Schur 维数乘积化简为
\((a+1)(b+1)(a+b+2)/2\)。这些等式说明分拆、最高权和张量构造是同一对象的不同描述。

\ProseParagraphBreak
后面的 BGG 方法研究 Verma 模之间的映射和消解，能够把特征标中的交错和解释为某些正合复形的 Euler 和。本章没有使用这项结构，也没有先假设相关消解存在。当前证明中的交错符号来自 Weyl 反对称性，额外项的消失来自 Casimir 的范数限制；这已经足够独立建立有限维特征标公式。
